\section{Limitations and evidence statement}\label{sec:scope}

The classification is exact for the specified ordered word and integer
domain. It gives a finite construction for every coefficient triple,
not a sharp universal table of periods or a complexity-efficient
implementation. The constants 27, 54 and \(100(H+1)\) are certified
bounds from the argument; no optimality is asserted. The equal-forcing
case, general Fricke geometry, group-orbit reduction and routine
finite-graph or zeta identities are not independent contributions of
this article.

The new ingredient is the all-level implication in
Proposition~\ref{prop:exhaustion}. It is proved by inequalities,
exact polynomial identities and bijectivity. No mathematical program,
finite-core enumeration or numerical experiment was executed for this
proof-only result. The pseudocode-style Procedure~\ref{proc:atlas}
states a terminating algorithm; it is not an execution receipt. The
previous equal-forcing manuscript's computer-assisted certificate is
credited but is neither imported as a proof dependency nor rerun.
There is no experimental dataset underlying the theorem.

The article was prepared with AI assistance. The underlying argument
received internal review within the current research team; this does
not represent external human peer review. Source comparisons are
bounded by the cited versions and inspected statements. The accompanying
source audit records access limits and bibliographic provenance;
unperformed retraction, venue-policy and conflict-of-interest checks
are not represented as completed checks.

Finally, the ordinary counts in \eqref{eq:fixed_counts} and
\eqref{eq:zeta} concern this source dynamical system only. They do not
identify target Euler factors, root numbers, an automorphic object,
or a Hilbert--P\'olya realization. No such arithmetic interpretation
is inferred from the existence of a finite cycle product.
